% Member 4 - final report, V Conclusion and Future Work
\section{Conclusion and Future Work}
\label{sec:conclusion}

We set out to find out what a style augmentation actually buys, and the answer turned out to be less than the method's premise suggests. Under a single-source Photo-only protocol on PACS, adding styled images to the training set improves cross-domain accuracy substantially, and the improvement survives the controls we designed to challenge it: doubling the number of training rows with byte-identical copies does not reproduce it, and neither does doubling them with non-stylistic transforms. The part of the gain that remains is attributable to the style itself.

The adaptive variant did not survive its control. A spatial strength map that applies less style where the content's structure is strong necessarily applies less style in total than a uniform mix at the same nominal strength, so a direct comparison confounds two factors. Once we hold the average strength fixed --- with a strength-matched uniform control and with a spatial map recalibrated to the same mean --- the two methods are indistinguishable in accuracy, and the per-image edge-retention analysis shows that the modulation does not preserve structure more selectively than the uniform alternative either. The mechanism the method was designed around is not observable in the images, and it is not observable in the classification result.

We draw three conclusions. First, style augmentation is a real and useful intervention on PACS, and it is worth the extra generation cost. Second, its benefit should be attributed to the amount of style applied rather than to the spatial structure preservation that motivated one line of work on it; the dose--response curve over $\alpha$ makes this explicit. Third, and most usefully for the field, the controls that make these attributions possible are cheap: a duplicate control and a strength-matched control are one extra generator and one extra training run each, and without them a headline gain can be reported for the wrong reason.

Several directions remain open. The soft Sobel response that drives the spatial map cannot tell a class-relevant contour from a texture edge; a semantic segmentation model would give a more selective map, and whether a selective map helps is the natural next experiment. The evidence here comes from one architecture and one dataset, so a ResNet-50 and a second domain-generalization benchmark would test how far the conclusion travels. And the third row of the decomposition remains an upper bound: separating "styling against a generic transform" from "one family of transforms against a random one" would need a finer set of controls. We are releasing the run artefacts and the checking machinery so that the decomposition can be re-run and extended.
